\documentclass[12pt,a4paper]{article}

% ============================================================
% PACKAGES
% ============================================================
\usepackage[utf8]{inputenc}
\usepackage[T1]{fontenc}
\usepackage[margin=2.5cm]{geometry}
\usepackage{array}
\usepackage{booktabs}
\usepackage{longtable}
\usepackage[table,dvipsnames]{xcolor}
\usepackage[hidelinks,colorlinks=true]{hyperref}
\usepackage{enumitem}
\usepackage{titlesec}
\usepackage{fancyhdr}
\usepackage[most]{tcolorbox}
\usepackage{tabularx}
\usepackage{setspace}
\usepackage{parskip}
\usepackage{graphicx}
\usepackage{float}

% ============================================================
% COLOR PALETTE
% ============================================================
\definecolor{accentdark}{HTML}{334155}
\definecolor{accentlight}{HTML}{DCE3EA}
\definecolor{principlebg}{HTML}{EEF1F4}
\definecolor{principleborder}{HTML}{334155}
\definecolor{grayrule}{HTML}{CCCCCC}
\definecolor{alertbg}{HTML}{F5ECEC}
\definecolor{alertborder}{HTML}{8C3A3A}
% heat-map: muted traffic light
\definecolor{heatlow}{HTML}{DCE6DD}
\definecolor{heatmed}{HTML}{EFE6D2}
\definecolor{heathigh}{HTML}{E6D2D2}

% ============================================================
% HYPERREF
% ============================================================
\hypersetup{
    linkcolor=accentdark,
    urlcolor=accentdark,
    citecolor=accentdark,
    pdftitle={GenAI-Assisted Document Redaction Governance},
    pdfauthor={Andrea Gennari}
}

% ============================================================
% SECTION STYLING
% ============================================================
\titleformat{\section}
  {\Large\bfseries\color{accentdark}}
  {\thesection}{1em}{}
  [\vspace{2pt}{\color{accentdark}\titlerule[0.7pt]}]
\titlespacing*{\section}{0pt}{24pt plus 4pt minus 2pt}{10pt plus 2pt minus 2pt}

\titleformat{\subsection}
  {\large\bfseries\color{accentdark!85!black}}
  {\thesubsection}{0.8em}{}
\titlespacing*{\subsection}{0pt}{18pt plus 3pt minus 2pt}{6pt plus 2pt minus 1pt}

% ============================================================
% HEADER / FOOTER
% ============================================================
\setlength{\headheight}{16pt}
\addtolength{\topmargin}{-4pt}
\pagestyle{fancy}
\fancyhf{}
\fancyhead[L]{\small\color{accentdark!70!black}\textit{GenAI-Assisted Document Redaction Governance}}
\fancyhead[R]{\small\color{accentdark!70!black}\textit{\nouppercase{\rightmark}}}
\fancyfoot[C]{\small\color{accentdark!70!black}\thepage}
\renewcommand{\headrulewidth}{0.4pt}
\renewcommand{\headrule}{\hbox to\headwidth{\color{accentdark!40!white}\leaders\hrule height \headrulewidth\hfill}}
\renewcommand{\footrulewidth}{0pt}

% ============================================================
% PARAGRAPH & LIST SPACING
% ============================================================
\setlength{\parskip}{6pt plus 1pt minus 1pt}
\setlength{\parindent}{0pt}

\setlist[itemize]{
  topsep=4pt, itemsep=3pt, parsep=1pt,
  leftmargin=1.8em,
  label=\textcolor{accentdark}{\textbullet}
}
\setlist[enumerate]{
  topsep=4pt, itemsep=3pt, parsep=1pt,
  leftmargin=2.2em,
  label=\textcolor{accentdark}{\arabic*.}
}

% ============================================================
% TCOLORBOX STYLES
% ============================================================
\newtcolorbox{definitionbox}[1][]{%
  enhanced, breakable, sharp corners,
  colback=white,
  colframe=accentdark,
  boxrule=0.5pt,
  coltitle=accentdark,
  colbacktitle=white,
  fonttitle=\bfseries\scshape,
  titlerule=0.4pt,
  toptitle=2pt, bottomtitle=2pt,
  left=8pt, right=8pt, top=6pt, bottom=6pt,
  title={#1},
}

\newtcolorbox{principlebox}[1][]{%
  enhanced, breakable, sharp corners,
  colback=principlebg,
  colframe=accentdark,
  boxrule=0.5pt,
  coltitle=accentdark,
  colbacktitle=principlebg,
  fonttitle=\bfseries\scshape,
  titlerule=0.4pt,
  borderline west={2pt}{0pt}{accentdark},
  toptitle=2pt, bottomtitle=2pt,
  left=10pt, right=8pt, top=6pt, bottom=6pt,
  title={#1},
}

\newtcolorbox{alertbox}[1][]{%
  enhanced, breakable, sharp corners,
  colback=white,
  colframe=alertborder,
  boxrule=0.6pt,
  coltitle=alertborder,
  colbacktitle=white,
  fonttitle=\bfseries\scshape,
  titlerule=0.4pt,
  toptitle=2pt, bottomtitle=2pt,
  left=8pt, right=8pt, top=6pt, bottom=6pt,
  title={#1},
}

\newtcolorbox{quotebox}{%
  breakable, sharp corners,
  colback=white,
  colframe=white,
  boxrule=0pt,
  borderline west={2pt}{0pt}{accentdark},
  left=10pt, right=8pt, top=4pt, bottom=4pt,
  fontupper=\itshape,
}

% ============================================================
% DOCUMENT BEGIN
% ============================================================
\begin{document}

\begin{titlepage}
    \centering
    \vspace*{3cm}

    {\color{accentdark}\rule{\textwidth}{1.5pt}}

    \vspace{1.2cm}
    {\Huge\bfseries\color{accentdark} GenAI-Assisted Document\par}
    \vspace{0.3cm}
    {\Huge\bfseries\color{accentdark} Redaction Governance\par}
    \vspace{0.8cm}
    {\Large\color{accentdark!80!black} A Governance Framework for Safe Adoption\par}
    \vspace{1.2cm}

    {\color{accentdark}\rule{0.4\textwidth}{0.8pt}}

    \vspace{2cm}
    {\large\color{accentdark!70!black} Submitted by\par}
    \vspace{0.3cm}
    {\LARGE\bfseries\color{accentdark} Andrea Gennari\par}
    \vspace{0.5cm}
    {\normalsize\color{accentdark!70!black} Business and Project Management\par}
    \vfill
    {\color{accentdark}\rule{\textwidth}{1.5pt}}
    \vspace{0.5cm}
    {\large\color{accentdark!70!black} \today\par}
\end{titlepage}

\thispagestyle{empty}
\vspace*{2cm}
\begin{tcolorbox}[
  enhanced, sharp corners,
  colback=white,
  colframe=accentdark,
  boxrule=0.5pt,
  coltitle=accentdark,
  colbacktitle=white,
  fonttitle=\bfseries\scshape,
  title=Abstract,
  titlerule=0.4pt,
  toptitle=2pt, bottomtitle=2pt,
  left=12pt, right=12pt, top=8pt, bottom=8pt,
]
This paper proposes a governance framework for the safe adoption of Generative AI-assisted document redaction tools. The central research question is: how can organizations adopt GenAI-assisted redaction through governance, human-in-the-loop controls, audit logging and risk management? The framework is organized in four pillars anchored to GDPR, the NIST AI Risk Management Framework and ISO/IEC 42001, and is illustrated against the Post Office Horizon data breach case.
\end{tcolorbox}

\newpage

\tableofcontents
\newpage

% ============================================================
% PART I — THE GOVERNANCE FRAMEWORK
% ============================================================
\begin{center}
{\color{accentdark}\rule{\textwidth}{1.2pt}}\\[6pt]
{\Huge\bfseries\color{accentdark} Part I}\\[8pt]
{\LARGE\bfseries\color{accentdark!85!black} A Governance Framework for Safe Adoption}\\[8pt]
{\color{accentdark}\rule{\textwidth}{1.2pt}}
\end{center}
\vspace{8pt}

\begin{quotebox}
\textbf{Guiding question for Part I.} How can organizations adopt GenAI-assisted document redaction safely, through governance, human-in-the-loop controls, audit logging and risk management?
\end{quotebox}

% ============================================================
\section{Framing and Foundations}
\label{sec:foundations}

\subsection{Framework Overview}

Organizations that publish documents externally need a controlled, auditable process to identify and redact sensitive content before release. GenAI can support this, but adopting it without governance adds new risks: over-reliance on automation, lost accountability, and data leakage.

This paper proposes a \textbf{governance framework} for that path, in four pillars, each addressing a distinct failure mode: Governance, Human-in-the-Loop Controls, Audit Logging, and Risk Management. Table~\ref{tab:structure} shows the structure.

\begin{table}[H]
\centering
\small
\rowcolors{2}{accentlight!30!white}{white}
\begin{tabular}{@{} p{4.8cm} p{9cm} @{}}
\toprule
\textbf{\color{accentdark}Section} & \textbf{\color{accentdark}Content} \\
\midrule
\S\ref{sec:foundations}: Foundations & Criteria for safe adoption and regulatory baseline \\
\S\ref{sec:governance}: Pillar~I & Governance, policy and accountability structure \\
\S\ref{sec:hitl}: Pillar~II & Human-in-the-loop controls and document risk classification \\
\S\ref{sec:audit}: Pillar~III & Audit logging and demonstrable compliance \\
\S\ref{sec:risk}: Pillar~IV & Risk management process and register \\
\S\ref{sec:adoption} & Adoption roadmap, KPIs and business case \\
\bottomrule
\end{tabular}
\caption{Structure of the paper.}
\label{tab:structure}
\end{table}

It is designed in full, then deployed through a limited pilot before scaling: predictive design, iterative rollout, a hybrid lifecycle \cite{falchi_slides}.

\subsection{Defining ``Safe Adoption''}
\label{sec:safe_adoption}

The framework must be verifiable. Five criteria define safe adoption; the framework is designed to satisfy all five at once.

\begin{definitionbox}[Five Criteria for Safe Adoption]
\begin{enumerate}
  \item \textbf{No unauthorized disclosure.} No personal or confidential data leaves without passing documented review and approval. (The Post Office Horizon incident is the motivating failure.)

  \item \textbf{Regulatory compliance.} Satisfies GDPR Art.~25 (protection by design) and Art.~32 (appropriate measures), and the ICO expectation of a documented pre-publication control.

  \item \textbf{Accountability and auditability.} Every document yields a record of what was done, by whom, and when, demonstrable to a regulator or court (GDPR Art.~5(2)).

  \item \textbf{Preserved human oversight.} No high-risk document is published without explicit human approval. GenAI \emph{supports}, it does not \emph{replace}, the reviewer, who is the final gate.

  \item \textbf{Sustainable adoption.} The framework is embedded: staff trained, process followed, KPIs met in the pilot. A correct framework that is bypassed reduces no risk.
\end{enumerate}
\end{definitionbox}

Each criterion has a failure condition that triggers escalation (Table~\ref{tab:failure}).

\begin{table}[H]
\centering
\small
\rowcolors{2}{accentlight!30!white}{white}
\begin{tabular}{@{} p{4.5cm} p{9cm} @{}}
\toprule
\textbf{\color{accentdark}Criterion} & \textbf{\color{accentdark}Failure condition (triggers escalation)} \\
\midrule
No unauthorized disclosure & Any unredacted or confidential content shared/published \\
Regulatory compliance & Any document processed outside the GDPR Art.~25/32 controls \\
Accountability \& auditability & Any document with an incomplete or missing audit record \\
Preserved human oversight & Any high-risk document released without documented human approval \\
Sustainable adoption & KPI targets missed beyond tolerance; systematic bypass of the process \\
\bottomrule
\end{tabular}
\caption{Failure conditions for each safe-adoption criterion.}
\label{tab:failure}
\end{table}

Each criterion is SMART: specific, measurable, attainable in the pilot, relevant, time-bound \cite{falchi_slides}.

\subsection{Regulatory and Standards Baseline}
\label{sec:baseline}

Three layers constrain any GenAI redaction tool: EU data-protection law, UK supervisory expectations, and international management-systems standards. Together they set the non-negotiable floor.

\noindent\textbf{GDPR: Key Articles}\par

\textbf{Art.~25 (by design and by default)}: protection must be built into the system \emph{before} processing, not retrofitted. Here: the review gate is part of the default publication path, not a skippable add-on.

\textbf{Art.~32 (security of processing)}: ``appropriate technical and organisational measures.'' Technical $=$ GenAI detection plus mandatory human review; organisational $=$ the policy, training, and audit log.

\textbf{Art.~5(2) (accountability)}: the controller must \emph{demonstrate} compliance. This drives Pillar~III: without a tamper-evident record of each decision, the organization cannot prove to the ICO or a court that the controls operated.

\noindent\textbf{ICO Expectations}\par

The ICO's 2025 reprimand of Post Office Limited held that publishing personal data without a documented pre-publication control breaches Art.~5(1)(f) and Art.~32 \cite{ico2025}. Its corrective requirement, a documented and auditable review process, is the origin of Pillars~II and~III, and it sets the bar: not just having a process, but proving it ran for every document.

\noindent\textbf{NIST AI Risk Management Framework (AI RMF 1.0)}\par

The NIST AI RMF \cite{nist_ai_rmf} organises AI risk governance around four functions:

\begin{itemize}
  \item \textbf{GOVERN}: AI policies, accountability, and risk culture.
  \item \textbf{MAP}: context, potential harms, and risk categorisation.
  \item \textbf{MEASURE}: quantify and evaluate risks with metrics and testing.
  \item \textbf{MANAGE}: prioritise, respond to, and monitor risks.
\end{itemize}

The framework's four pillars map directly: Pillar~I (Governance) operationalises GOVERN; Pillar~IV (Risk Management) covers MAP, MEASURE, and MANAGE.

\noindent\textbf{ISO/IEC Standards}\par

\textbf{ISO/IEC 42001:2023} \cite{iso42001} (AI management): pre-deployment risk assessment (Cl.~6), documented human oversight (A.6), impact evaluation on data subjects (A.9). Met by Pillars~II, III, and~IV.

\textbf{ISO/IEC 27001:2022} \cite{iso27001} (information security): information classification (A.8.2), data masking (A.8.11), security policy (A.5.1). Pillar~II implements A.8.2; the audit log meets A.8.15.

\noindent\textbf{Requirement Mapping}\par

Table~\ref{tab:standards} maps each requirement to its source and pillar; every row ties back to a safe-adoption criterion (Section~\ref{sec:safe_adoption}).

\begin{table}[H]
\centering
\small
\rowcolors{2}{accentlight!30!white}{white}
\begin{tabular}{@{} p{4.3cm} p{3.5cm} p{5.7cm} @{}}
\toprule
\textbf{\color{accentdark}Requirement} & \textbf{\color{accentdark}Source} & \textbf{\color{accentdark}Addressed in} \\
\midrule
Data protection by design & GDPR Art.~25 & Pillar~I (policy mandate) + Pillar~II (default review gate) \\
Appropriate security measures & GDPR Art.~32 & Pillar~I + II + III \\
Demonstrable accountability & GDPR Art.~5(2) & Pillar~III (immutable audit log) \\
Documented pre-publication control & ICO reprimand 2025 & Pillar~II (HITL review gate) \\
AI governance and risk culture & NIST AI RMF GOVERN & Pillar~I (governance structure, RACI) \\
AI risk identification & NIST AI RMF MAP & Pillar~IV (risk register) \\
AI risk measurement & NIST AI RMF MEASURE & Pillar~IV (performance metrics and KPIs) \\
AI risk treatment & NIST AI RMF MANAGE & Pillar~IV (response strategies) \\
AI human oversight & ISO/IEC 42001 A.6 & Pillar~II (mandatory human approval) \\
AI impact assessment & ISO/IEC 42001 Cl.~6 & Pillar~IV (pre-deployment risk assessment) \\
Information classification & ISO/IEC 27001 A.8.2 & Pillar~II (document risk tiers) \\
Audit evidence & ISO/IEC 27001 A.8.15 & Pillar~III (audit log) \\
\bottomrule
\end{tabular}
\caption{Regulatory and standards requirements mapped to framework pillars.}
\label{tab:standards}
\end{table}

\subsection{Framework Architecture}
\label{sec:architecture}

The framework is a \textbf{governance spine} (Pillar~I) giving authority to three operational pillars, all feeding a hybrid adoption lifecycle. One design principle governs everything that follows:

\begin{principlebox}[Architectural Principle]
\textbf{AI suggests, humans approve, the organization remains accountable.} The tool is a detection assistant, not a decision authority: every publication decision is owned by a named human and recorded in the log, and the policy, not the tool's confidence score, sets when approval is mandatory.
\end{principlebox}

\noindent\textbf{The Governance Spine}\par

Pillar~I is load-bearing: the policy mandate, the accountability hierarchy, and the RACI matrix that authorize the other three pillars. Without it the log is not required, the gate can be bypassed, and the register has no owner.

\noindent\textbf{The Three Operational Pillars}\par

\begin{itemize}
  \item \textbf{Pillar~II: Human-in-the-Loop Controls} classifies every document into three risk tiers (standard, high, very high) and routes it through the matching review gate. The floor is \emph{standard}, not \emph{low}: every document holds personal data, so none skips review. High and very-high documents need explicit human approval; GenAI output is a first pass, not a decision.

  \item \textbf{Pillar~III: Audit Logging} records every action on every document (proposals, each review step, the final decision, plus actor and timestamp) in an immutable, queryable log retained as regulatory evidence.

  \item \textbf{Pillar~IV: Risk Management} keeps a living risk register for the tool, tracks residual risk after controls, triggers escalation at thresholds, and defines the manual fallback for outages.
\end{itemize}

\noindent\textbf{Document Flow Through the Framework}\par

Table~\ref{tab:docflow} shows the path from draft to publication: linear for standard documents, with high and very-high adding the approval step (Step~5) before the token is issued.

\begin{table}[H]
\centering
\small
\rowcolors{2}{accentlight!30!white}{white}
\begin{tabular}{@{} c p{3.2cm} p{8.3cm} @{}}
\toprule
\textbf{\color{accentdark}Step} & \textbf{\color{accentdark}Actor} & \textbf{\color{accentdark}Action} \\
\midrule
1 & Document owner & Submits draft to the redaction workflow \\
2 & Pillar~II system & Classifies document risk tier; routes to appropriate review queue \\
3 & GenAI tool & Proposes redactions; confidence scores recorded in audit log (Pillar~III) \\
4 & Human reviewer & Reviews, confirms or overrides GenAI proposals; decision recorded \\
5 & Approval authority & \textit{(High/very-high risk only)} Signs off on final redacted version \\
6 & Pillar~III log & Locks the complete review record; approval token generated \\
7 & Document owner & Publishes against the approval token \\
\bottomrule
\end{tabular}
\caption{Document flow from draft to publication under the framework.}
\label{tab:docflow}
\end{table}

\noindent\textbf{Hybrid Adoption Lifecycle}\par

The framework is designed predictively (all four pillars defined before the pilot) but deployed iteratively: a controlled pilot generates KPI data and calibrates the risk-tier thresholds and review workloads before full rollout. This hybrid sequencing (Section~\ref{sec:foundations}) cuts adoption risk without losing coherence \cite{falchi_slides}.

% ============================================================
\section{Governance (Pillar I)}
\label{sec:governance}

Pillar~I is the governance spine (Section~\ref{sec:architecture}): it gives the other three pillars their authority through four artefacts: an operating model, a redaction policy, decision rights, and a RACI matrix. Without it the audit log is optional, the review gate is skippable, and the risk register has no owner.

\subsection{Operating Model and Oversight}

The framework is owned end to end by the \textbf{Data Protection Officer (DPO)}, who chairs a \textbf{Redaction Governance Board}: the oversight forum that approves the policy, sets the Pillar~II risk thresholds, reviews KPI and audit data (Section~\ref{sec:audit}), and authorizes exceptions. Five roles staff the workflow:

\begin{definitionbox}[Governance Roles]
\begin{itemize}
  \item \textbf{Document Owner}: submits the document, runs the GenAI scan, performs first-line review, and publishes against the approval token.
  \item \textbf{Legal}: advises on litigation, contractual and disclosure sensitivity; recommends on high-risk release.
  \item \textbf{DPO / Privacy}: accountable owner of the framework; chairs the Board; holds classification and high-risk approval authority; attests to audit-log integrity.
  \item \textbf{IT Security}: owns the GenAI pipeline, the vendor assessment, access controls and the log infrastructure.
  \item \textbf{Manager}: holds operational accountability for the routine flow; provides resources and owns the KPIs.
\end{itemize}
\end{definitionbox}

These parties differ in power and interest, so the framework engages them differently \cite{falchi_slides}:

\begin{itemize}
  \item \textbf{DPO, Legal} (high power, high interest): \emph{manage closely}; on the Board, with approval rights.
  \item \textbf{IT Security} (high power, lower interest): \emph{keep satisfied}; owns the controls, consulted per document.
  \item \textbf{Document owners} (lower power, high interest): \emph{keep informed} and trained.
  \item \textbf{Senior management} (high power): \emph{keep satisfied} via the Board's reporting line.
\end{itemize}

Communication splits the same way: approval decisions are \emph{pull} (recorded, retrieved on demand); policy changes and incidents are \emph{push} to every owner \cite{falchi_slides}.

\subsection{GenAI-Assisted Redaction Policy}

The policy is the artefact GDPR Art.~25 demands: it embeds data protection into the process \emph{by default}, rather than leaving it to individual diligence under time pressure.

\begin{definitionbox}[Redaction Policy at a Glance]
\begin{itemize}
  \item \textbf{Purpose}: ensure no document is shared or published externally without a documented, risk-proportionate redaction review.
  \item \textbf{Scope}: every document leaving the organization to a named external party or the public, and all staff who publish.
  \item \textbf{Core principle}: AI suggests, humans approve, the organization remains accountable (Section~\ref{sec:architecture}).
  \item \textbf{Prohibited uses}: no automatic publication at any tier; no pasting of documents into consumer GenAI tools outside the approved pipeline; no disabling of the audit log.
  \item \textbf{Vendor \& data handling}: a Data Processing Agreement (GDPR Art.~28), a contractual no-training clause, EU or approved data residency, and encryption in transit and at rest.
\end{itemize}
\end{definitionbox}

The vendor rules close a risk nothing else can: sending text to an external model is itself a disclosure. They are the control for the ``data exposure to external AI vendor'' risk (Section~\ref{sec:risk}).

\subsection{Decision Rights and Approval Gates}

Decision rights answer one question (\emph{who may authorize external release?}); the answer rises with the risk tier (Section~\ref{sec:hitl}). No tier permits automatic release: a human must record approval before a document advances.

\begin{table}[H]
\centering
\small
\rowcolors{2}{accentlight!30!white}{white}
\begin{tabular}{@{} p{2.6cm} p{4.8cm} p{6.4cm} @{}}
\toprule
\textbf{\color{accentdark}Risk tier} & \textbf{\color{accentdark}Release authority} & \textbf{\color{accentdark}Gate rule} \\
\midrule
Standard & Document owner & Owner attests review complete; token issued on attestation \\
High & DPO, on Legal's recommendation & No release until DPO approval is recorded; the author cannot self-approve \\
Very high & DPO \emph{and} senior management (dual control) & Two recorded approvals required; no single person can release \\
\bottomrule
\end{tabular}
\caption{Decision rights and approval gates by risk tier.}
\label{tab:decisionrights}
\end{table}

A \textbf{gate} is a hard stop: no publication until the named authority records approval in the log (Pillar~III). The very-high tier adds \emph{dual control} so no single person can release the most sensitive material. These gates operationalize the ``preserved human oversight'' criterion (Section~\ref{sec:safe_adoption}).

\subsection{RACI Matrix}

Every activity has exactly one \textbf{Accountable} and at least one \textbf{Responsible} role, with Consulted and Informed named alongside \cite{falchi_slides}. Operational steps answer to the Manager; privacy-critical control points answer to the DPO.

\begin{table}[H]
\centering
\small
\rowcolors{2}{accentlight!30!white}{white}
\begin{tabular}{@{} >{\bfseries\color{accentdark}}p{4.4cm} *{5}{>{\centering\arraybackslash}p{1.75cm}} @{}}
\toprule
\textbf{\color{accentdark}Activity} & \textbf{\color{accentdark}Doc Owner} & \textbf{\color{accentdark}Legal} & \textbf{\color{accentdark}DPO} & \textbf{\color{accentdark}IT Sec.} & \textbf{\color{accentdark}Manager} \\
\midrule
Intake \& version check & R & I & I & C & A \\
GenAI sensitive-data scan & R & I & C & A & I \\
Risk-tier classification & R & C & A & C & I \\
Redaction proposal review & R & C & C & I & A \\
High-risk release approval & I & R & A & C & C \\
Final publication & R & I & I & C & A \\
Audit review & I & C & A & R & I \\
\bottomrule
\end{tabular}
\caption{RACI matrix for the redaction governance workflow.}
\label{tab:raci}
\end{table}

\begin{tcolorbox}[
  enhanced, sharp corners, colback=white, colframe=gray!45, boxrule=0.5pt,
  left=8pt, right=8pt, top=5pt, bottom=5pt,
]
\textbf{Legend:}\quad R\,=\,Responsible\quad A\,=\,Accountable\quad C\,=\,Consulted\quad I\,=\,Informed
\end{tcolorbox}

Two assignments carry the weight. \textbf{High-risk approval}: Accountable to the DPO, Responsible to Legal, so it is single-owned yet never decided without an independent recommendation. \textbf{Audit review}: Responsible to IT Security (who runs the log), Accountable to the DPO (who attests to it), preserving separation of duties. Every Accountable cell is a Board member.

% ============================================================
\section{Human-in-the-Loop Controls (Pillar II)}
\label{sec:hitl}

Pillar~II turns ``AI suggests, humans approve'' into procedure: it classifies every document into a risk tier, binds each tier to a level of human control, and routes it through gated review and escalation. Human attention is spent in proportion to risk.

\subsection{Document Classification and Risk Scoring}

Classification is the first control and runs on \emph{every} document. Three factors are scored; their product sets the tier.

\begin{definitionbox}[Risk Score $=$ Sensitivity $\times$ Exposure $\times$ Audience]
\begin{itemize}
  \item \textbf{Sensitivity}: the data categories present: 1~$=$ routine business; 2~$=$ personal data; 3~$=$ special-category, financial or litigation data.
  \item \textbf{Exposure}: the number of data subjects affected: 1~$=$ few; 2~$=$ tens; 3~$=$ hundreds or more.
  \item \textbf{Audience}: the reach on release: 1~$=$ internal; 2~$=$ a named external party; 3~$=$ the public web.
\end{itemize}
\end{definitionbox}

The score (1 to 27) maps to the three tiers of Section~\ref{sec:architecture}: \textbf{standard} ($\leq 4$), \textbf{high} ($6$ to $12$), \textbf{very high} ($\geq 18$). Override: any \emph{public} document (Audience~$=3$) is floored at \textbf{high}, since public release is irreversible. The Post Office document scored the maximum ($3\times3\times3=27$); Part~II shows the framework would have flagged it very high.

\subsection{Risk-Tier to Control Matrix}

The tier fixes the level of human control. Across every tier the GenAI tool only \emph{proposes}; it never publishes. What rises with risk is the seniority and number of humans who must sign off.

\begin{table}[H]
\centering
\small
\rowcolors{2}{accentlight!30!white}{white}
\begin{tabular}{@{} p{2.1cm} p{3.6cm} p{4.7cm} p{3.4cm} @{}}
\toprule
\textbf{\color{accentdark}Tier} & \textbf{\color{accentdark}GenAI role} & \textbf{\color{accentdark}Mandatory human control} & \textbf{\color{accentdark}Release authority} \\
\midrule
Standard & Proposes redactions; flags candidates & Document owner reviews every proposal and attests completeness & Document owner \\
High & Proposes redactions & Independent reviewer (not the author) plus compliance check & DPO, on Legal's recommendation \\
Very high & Proposes redactions & Independent review plus dual approval; no single-person release & DPO + senior management \\
\bottomrule
\end{tabular}
\caption{Human control required at each risk tier. No tier allows automatic release.}
\label{tab:risktier}
\end{table}

The lowest tier is \emph{standard}, not \emph{low}: every document is reviewed, none is exempt.

\subsection{Review and Escalation Workflow}

The workflow adds to the flow of Table~\ref{tab:docflow}, at each step, the control gate (what must hold before advancing) and the escalation branch.

\begin{table}[H]
\centering
\small
\rowcolors{2}{accentlight!30!white}{white}
\begin{tabular}{@{} c p{2.9cm} p{4.1cm} p{4.9cm} @{}}
\toprule
\textbf{\color{accentdark}\#} & \textbf{\color{accentdark}Actor (lane)} & \textbf{\color{accentdark}Action} & \textbf{\color{accentdark}Gate / escalation} \\
\midrule
1 & Document owner & Submit draft and intended audience & (none) \\
2 & System & Score and classify the tier & Public audience floors the tier at high \\
3 & GenAI tool & Propose redactions with confidence scores & Low-confidence items flagged for mandatory attention \\
4 & Reviewer & Confirm or override each proposal; attest completeness & High / very high require an \emph{independent} reviewer \\
5 & Approver & Approve the redacted version & High: DPO gate. Very high: DPO + senior (dual). A disagreement between reviewer and approver escalates to the Board \\
6 & Audit log & Lock the record; issue approval token & No token without a complete record \\
7 & Document owner & Publish against the token & Publication blocked without a valid token \\
\bottomrule
\end{tabular}
\caption{Gated review-and-escalation workflow (swimlane by actor).}
\label{tab:workflow}
\end{table}

\subsection{Reviewer Guidance and Automation Bias}

The weak point of any human-in-the-loop design is \emph{automation bias}: if reviewers rubber-stamp the tool, the gate is theatre and the over-reliance risk (Section~\ref{sec:risk}) materializes. Four choices counter it:

\begin{itemize}
  \item \textbf{Show the original, not just the result.} The reviewer sees the source document beside the proposals, so a \emph{missing} redaction is visible rather than hidden in clean-looking output.
  \item \textbf{Make confirmation an act, not a default.} The reviewer positively confirms each proposed redaction, and a ``no sensitive data'' outcome requires an explicit attestation: a recorded decision, never a silent pass.
  \item \textbf{Measure catch rate.} Blind quality samples (documents seeded with known sensitive items) test whether reviewers catch what the tool misses; the catch rate is a KPI (Section~\ref{sec:adoption}).
  \item \textbf{Gate on competence.} Only trained reviewers are provisioned, and high-risk review requires additional training, the control for the ``lack of staff training'' risk.
\end{itemize}

In PM terms, classification is probability/impact prioritization and escalation is structured decision-making \cite{falchi_slides}.

% ============================================================
\section{Audit Logging and Accountability (Pillar III)}
\label{sec:audit}

Pillar~III turns the workflow into evidence. GDPR Art.~5(2) requires the organization to \emph{demonstrate} compliance, not just achieve it, and the ICO's reprimand of Post Office Limited turned on the absence of a documented pre-publication process \cite{ico2025}. The log is the artefact that proves, for any document, that the Art.~32 controls operated.

\subsection{Purpose and Principles}

\begin{principlebox}[Three Properties of the Log]
\textbf{Complete}: every document that enters the workflow produces a record; no record, no publication token.\par
\textbf{Immutable}: entries are append-only and tamper-evident (hash-chained), so a record cannot be edited after the fact to look compliant.\par
\textbf{Traceable}: every entry binds an action to an identified actor and a timestamp, so the whole decision chain can be reconstructed.
\end{principlebox}

Completeness is enforced by the Pillar~II token gate; immutability earns a regulator's trust; traceability answers ``who decided this, and when?''

\subsection{Audit Log Schema}

Four field groups, plus the controls that protect the log itself:

\begin{definitionbox}[Audit Log Schema]
\begin{itemize}
  \item \textbf{Document identity}: document ID; type; version hash; owner; intended audience; classification tier.
  \item \textbf{Detection}: GenAI model and version; sensitive-data categories detected; per-item confidence scores; scan timestamp.
  \item \textbf{Decision}: reviewer identity; each proposal confirmed or overridden; changes made; completeness attestation; escalation flag.
  \item \textbf{Disposition}: approver identity; approval or rejection; approval timestamp; approval token ID; publication status.
\end{itemize}
\textbf{Controls on the log:} append-only (no edit or delete for any role); read access restricted to DPO and IT Security; access to the log is itself logged; retention aligned to the limitation period for related claims, then deletion (GDPR Art.~5(1)(e), storage limitation).
\end{definitionbox}

Recording the model version and per-item confidence lets an investigation separate a tool failure (item never flagged) from a reviewer failure (flagged but dismissed).

\subsection{Demonstrable Compliance}

The test of the log is whether it answers a regulator: to ``show that document~X was handled correctly,'' one query returns who scanned it, what was detected, who reviewed and changed it, who approved it, and when. Table~\ref{tab:compliance} maps log elements to requirements.

\begin{table}[H]
\centering
\small
\rowcolors{2}{accentlight!30!white}{white}
\begin{tabular}{@{} p{6.0cm} p{8.0cm} @{}}
\toprule
\textbf{\color{accentdark}Log element} & \textbf{\color{accentdark}Requirement satisfied} \\
\midrule
Complete per-document record & GDPR Art.~5(2) accountability; ICO documented-process expectation \\
Reviewer and approver identity + timestamp & Art.~32 (which control operated, and by whom); preserved-oversight criterion \\
Detected categories and redaction changes & Art.~32 (evidence the measure worked); ``no unauthorized disclosure'' criterion \\
Append-only, tamper-evident storage & Evidential integrity for the ICO or a court \\
Retention limit and restricted access & Art.~5(1)(e) storage limitation; Art.~32 confidentiality \\
\bottomrule
\end{tabular}
\caption{Audit-log elements mapped to the requirements they evidence.}
\label{tab:compliance}
\end{table}

\subsection{Audit Review Process}

The log is reviewed, not just retained: monthly, the DPO reviews a sample plus \emph{every} very-high record; quarterly, the Board reviews KPIs and exceptions. Any anomaly (missing approver, bypassed gate, incomplete record, KPI breach) is logged as an incident and escalated, closing the loop with Pillar~IV.

% ============================================================
\section{Risk Management (Pillar IV)}
\label{sec:risk}

Pillar~IV manages the risk of the tool itself. Automation removes some risk (human inconsistency, volume misses) and adds new risk (automation bias, vendor exposure, false confidence). The PMI process identifies, scores, treats, and monitors it \cite{falchi_slides}.

\subsection{Risk Management Approach}

The PMI process has six steps, and each already has a home in the framework rather than living as a separate exercise:

\begin{itemize}
  \item \textbf{Plan}: the redaction policy and this pillar's method (Pillar~I).
  \item \textbf{Identify}: the risk register below.
  \item \textbf{Qualitative analysis}: the probability/impact heat map.
  \item \textbf{Quantitative analysis}: expected monetary value for the costliest risks.
  \item \textbf{Plan responses}: one of avoid / transfer / mitigate / accept per risk, mapped to Pillars I, II and III.
  \item \textbf{Monitor \& control}: the audit review (Pillar~III) and KPIs (Section~\ref{sec:adoption}), with defined triggers.
\end{itemize}

\subsection{Risk Register}

Each risk carries an ID, category, and root cause, so responses target causes. P and I are \emph{inherent} ratings (pre-control); residuals are in Table~\ref{tab:responses}.

\begin{table}[H]
\centering
\small
\rowcolors{2}{accentlight!30!white}{white}
\begin{tabular}{@{} c p{2.2cm} p{3.8cm} p{3.8cm} c c @{}}
\toprule
\textbf{\color{accentdark}ID} & \textbf{\color{accentdark}Category} & \textbf{\color{accentdark}Risk} & \textbf{\color{accentdark}Root cause} & \textbf{\color{accentdark}P} & \textbf{\color{accentdark}I} \\
\midrule
R1 & Detection & Sensitive data not detected (false negative) & Model limits; novel data formats & M & H \\
R2 & Process & Wrong document version published & No version control; manual handling & M & H \\
R3 & Human factor & Over-reliance on AI (automation bias) & Reviewer defers to the tool & M & H \\
R4 & Compliance & Cannot demonstrate compliance & No persistent record of decisions & M & H \\
R5 & Vendor / data & Document text exposed to external AI vendor & Text sent off-premise for inference & M & H \\
R6 & Quality & Excess false positives erode trust and throughput & Conservative model tuning & H & M \\
R7 & Organizational & Staff untrained; process bypassed & Rushed rollout & M & M \\
R8 & Governance & Unclear accountability for release & No defined decision rights & M & H \\
\bottomrule
\end{tabular}
\caption{Risk register (inherent probability P and impact I, before controls).}
\label{tab:riskregister}
\end{table}

\subsection{Qualitative and Quantitative Analysis}

The probability/impact grid \cite{falchi_slides} shows the profile (Table~\ref{tab:heatmap}). The cluster in the medium-probability / high-impact cell justifies a control-heavy framework: these risks are neither rare nor mild.

\begin{table}[H]
\centering
\small
\rowcolors{2}{white}{white}
\begin{tabular}{@{} c >{\centering\arraybackslash}p{2.3cm} >{\centering\arraybackslash}p{2.3cm} >{\centering\arraybackslash}p{3.6cm} @{}}
\toprule
 & \multicolumn{3}{c}{\textbf{\color{accentdark}Impact} $\rightarrow$} \\
\textbf{\color{accentdark}Probability $\downarrow$} & \textbf{Low} & \textbf{Medium} & \textbf{High} \\
\midrule
\textbf{High}   & \cellcolor{heatmed} & \cellcolor{heathigh} R6 & \cellcolor{heathigh} \\
\textbf{Medium} & \cellcolor{heatlow} & \cellcolor{heatmed} R7 & \cellcolor{heathigh} R1, R2, R3, R4, R5, R8 \\
\textbf{Low}    & \cellcolor{heatlow} & \cellcolor{heatlow} & \cellcolor{heatmed} \\
\bottomrule
\end{tabular}
\caption{Probability/impact heat map of the inherent risk register.}
\label{tab:heatmap}
\end{table}

The costliest risks are weighed by \textbf{expected monetary value} (EMV $=$ probability $\times$ impact) \cite{falchi_slides}: a control pays off when the expected loss it removes exceeds its cost. Cutting a risk from 20\% to 2\% on a EUR 500,000 impact saves EUR 90,000 a year, far above the tool's cost. Two caveats: EMV is an average and hides the tail (the flaw of averages), so the very-high tier uses dual control rather than a cost-benefit test \cite{falchi_slides}; and these figures are illustrative, the case-calibrated EMV being deferred to the Part~II counterfactual (Section~\ref{sec:rq2_counterfactual}).

\subsection{Risk Responses and Residual Risk}

Each risk takes one of the four PMI strategies (avoid, transfer, mitigate, accept), delivered by a control in another pillar \cite{falchi_slides}: the framework is the response plan for its own register.

\begin{table}[H]
\centering
\small
\rowcolors{2}{accentlight!30!white}{white}
\begin{tabular}{@{} c p{2.4cm} p{8.2cm} >{\centering\arraybackslash}p{1.9cm} @{}}
\toprule
\textbf{\color{accentdark}ID} & \textbf{\color{accentdark}Strategy} & \textbf{\color{accentdark}Response (control $\rightarrow$ pillar)} & \textbf{\color{accentdark}Residual} \\
\midrule
R1 & Mitigate & Mandatory review of all proposals; independent review for high risk; catch-rate KPI ($\rightarrow$ II) & Low/Med \\
R2 & Mitigate & Version hash and token-gated publication ($\rightarrow$ II/III) & Low \\
R3 & Mitigate & Confirm-each-item review, original shown, blind samples, training ($\rightarrow$ II) & Low/Med \\
R4 & Mitigate & Immutable, complete audit log ($\rightarrow$ III) & Low \\
R5 & Transfer + mitigate & DPA and no-training clause shift liability; EU residency limits exposure ($\rightarrow$ I) & Low/Med \\
R6 & Accept (reduce) & Reviewer adjusts proposals; thresholds tuned in the pilot; residual false positives accepted as the cost of safety ($\rightarrow$ II) & Med \\
R7 & Avoid + mitigate & No provisioning without training; gated rollout ($\rightarrow$ I) & Low \\
R8 & Mitigate & RACI, decision rights and approval gates ($\rightarrow$ I) & Low \\
\bottomrule
\end{tabular}
\caption{Risk responses mapped to controls, with residual risk.}
\label{tab:responses}
\end{table}

All four strategies appear: most risks are \emph{mitigated}; vendor liability is partly \emph{transferred} by contract (R5); untrained-staff exposure is \emph{avoided} (R7); irreducible false positives are \emph{accepted} (R6). Residual risk is monitored via the audit review and KPIs, with \textbf{triggers} (catch rate below target, or any released-without-approval event) escalating to the Board and a contingency reserve for incidents \cite{falchi_slides}. The claim is not zero risk but controlled risk: every inherent risk moved out of the red cell of Table~\ref{tab:heatmap} and kept there.

% ============================================================
\section{Adoption, Measurement and Business Case}
\label{sec:adoption}
\textit{[To be developed.]}

% ============================================================
\section{Synthesis: The Framework for Safe Adoption}
\label{sec:synthesis}
\textit{[To be developed.]}

% ============================================================
% PART II — THE POST OFFICE HORIZON CASE
% ============================================================
\clearpage
\begin{center}
{\color{accentdark}\rule{\textwidth}{1.2pt}}\\[6pt]
{\Huge\bfseries\color{accentdark} Part II}\\[8pt]
{\LARGE\bfseries\color{accentdark!85!black} The Post Office Horizon Case}\\[8pt]
{\color{accentdark}\rule{\textwidth}{1.2pt}}
\end{center}
\vspace{8pt}

\begin{quotebox}
\textbf{Guiding question for Part II.} How could a GenAI-assisted governance framework have reduced the privacy and data leakage risks in the Post Office Horizon document publication case?
\end{quotebox}

% ============================================================
\section{Case Analysis: Root Causes and Governance Gaps}
\label{sec:rq2_case}
\textit{[To be developed.]}

% ============================================================
\section{Applying the Framework to the Case}
\label{sec:rq2_apply}

\subsection{Governance Controls (Pillar I)}
\textit{[To be developed.]}

\subsection{Human-in-the-Loop Controls (Pillar II)}
\textit{[To be developed.]}

\subsection{Audit Logging (Pillar III)}
\textit{[To be developed.]}

\subsection{Risk Management (Pillar IV)}
\textit{[To be developed.]}

% ============================================================
\section{Counterfactual Analysis}
\label{sec:rq2_counterfactual}

% What would have changed if the framework had been in place?
\textit{[To be developed: for each root cause identified in \S\ref{sec:rq2_case}, show which framework control would have prevented or mitigated it.]}

% ============================================================
\section{Synthesis: Lessons from the Case}
\label{sec:rq2_synthesis}
\textit{[To be developed.]}

% ============================================================
% REFERENCES
% ============================================================
\begin{thebibliography}{9}

\bibitem{falchi_slides}
Alessio Falchi.
\textit{Project \& Business Management: Course Slides}.
University of Pisa, 2024.

\bibitem{gdpr}
European Union.
\textit{General Data Protection Regulation (GDPR), Articles 5, 25 and 32}.

\bibitem{nist_ai_rmf}
National Institute of Standards and Technology.
\textit{Artificial Intelligence Risk Management Framework (AI RMF 1.0)}.
2023.

\bibitem{iso42001}
International Organization for Standardization.
\textit{ISO/IEC 42001:2023, Artificial Intelligence Management Systems}.
2023.

\bibitem{iso27001}
International Organization for Standardization.
\textit{ISO/IEC 27001:2022, Information Security Management Systems}.
2022.

\bibitem{ico2025}
Information Commissioner's Office.
\textit{Post Office Limited Reprimand}.
2025.

\end{thebibliography}

\end{document}
